% Standalone, self-contained LaTeX manuscript.
% Compile: pdflatex -interaction=nonstopmode Strict_Subadditivity_Fourth_Order_Cubic_Quintic.tex (twice)
\documentclass[11pt,letterpaper]{article}

\usepackage[T1]{fontenc}
\usepackage[utf8]{inputenc}
\usepackage{lmodern}
\usepackage{microtype}
\usepackage[margin=1.05in]{geometry}
\usepackage{amsmath,amssymb,amsthm,mathtools}
\usepackage{hyperref}
\usepackage{enumitem}
\usepackage{xcolor}

\hypersetup{
  colorlinks=true,
  linkcolor=blue!55!black,
  citecolor=blue!55!black,
  urlcolor=blue!55!black,
  pdftitle={Strict Subadditivity for a Fourth-Order Cubic--Quintic Hamiltonian in Two Dimensions},
  pdfauthor={Richard J. Reyes},
  pdfsubject={Mass-constrained fourth-order variational problem; explicit strict-binding threshold}
}
\setlength{\parindent}{1.15em}
\setlength{\parskip}{0.1em}
\allowdisplaybreaks[2]
\numberwithin{equation}{section}

\newtheorem{theorem}{Theorem}[section]
\newtheorem{lemma}[theorem]{Lemma}
\newtheorem{proposition}[theorem]{Proposition}
\newtheorem{corollary}[theorem]{Corollary}
\newtheorem{remark}[theorem]{Remark}
\newcommand{\RR}{\mathbb R}
\newcommand{\CC}{\mathbb C}
\newcommand{\norm}[1]{\left\lVert #1\right\rVert}
\newcommand{\abs}[1]{\left\lvert #1\right\rvert}
\newcommand{\dd}{\,\mathrm d}
\newcommand{\esssup}{\operatorname*{ess\,sup}}

\title{\bfseries Strict Subadditivity for a Fourth-Order\\
Cubic--Quintic Hamiltonian in Two Dimensions}
\author{Richard J. Reyes}
\date{October 26, 2026}

\begin{document}
\maketitle

\begin{abstract}
We consider the mass-constrained infimum of a fourth-order Hamiltonian with a focusing quartic energy term and a defocusing sextic energy term on $H^2(\RR^2;\CC)$. We establish strict subadditivity for every interior splitting of each mass $0<m\le m_B$, where
\[
  m_B=\frac{3\beta}{\gamma\sqrt{k^4+3\beta^2/(32\gamma)}}
  \qquad (\beta,\gamma,k>0).
\]
The proof first constructs negative-energy trials at every positive mass, then uses an exact mass-preserving dilation to replace arbitrary minimizing sequences by sequences satisfying a scalar stationarity identity. A Fourier $L^\infty$ estimate and a square-completion identity control the sextic-to-quartic ratio and yield a quantitative strict decrease of energy per unit mass. No attainment of the constrained infimum is assumed. An exact-arithmetic corollary verifies the sufficient condition for a specified frozen parameter set. The result is confined to this two-dimensional combined-power variational problem and does not assert a sharp critical mass or conclusions about physical confinement lifetimes.
\end{abstract}

\noindent\textbf{Keywords.} Fourth-order nonlinear Schr\"odinger energy; strict subadditivity; mass constraint; combined nonlinearities; concentration--compactness; mass-preserving dilation.\\
\textbf{Mathematics Subject Classification (2020).} 35Q55, 35J30, 49J40.

\section{Introduction and main result}

Let
\[
  X=H^2(\RR^2;\CC),\qquad A=\Delta+k^2,
\]
and, for $\beta,\gamma,k>0$, define
\begin{equation}\label{eq:H}
  H(u)=\norm{Au}_2^2-\frac{\beta}{2}\norm{u}_4^4
                +\frac{\gamma}{3}\norm{u}_6^6.
\end{equation}
For $m>0$, let
\begin{equation}\label{eq:I}
  I(m)=\inf\bigl\{H(u):u\in X,\ \norm{u}_2^2=m\bigr\}.
\end{equation}
All integrals are over $\RR^2$ with Lebesgue measure. No radial, real-valued, periodic, band-limited, or other symmetry restriction is imposed.

For orientation, mixed-dispersion fourth-order nonlinear Schr\"odinger variational problems have been studied by Fern\'andez, Jeanjean, Mandel, and Mari\c{s} \cite{FJMM2022}, and combined-power mixed-dispersion problems by Ma, Chang, and Feng \cite{MCF2024}. The result proved here is the explicit sufficient mass interval for strict subadditivity of the \emph{particular} Hamiltonian \eqref{eq:H}. No assertion of an exhaustive literature-based priority or optimality is made.

\begin{theorem}[Explicit strict-binding range]\label{thm:main}
Let $\beta,\gamma,k>0$ and set
\begin{equation}\label{eq:mB}
 \rho=k^4+\frac{3\beta^2}{32\gamma},\qquad
 \sigma=\sqrt{\rho},\qquad
 m_B=\frac{3\beta}{\gamma\sigma}.
\end{equation}
For every $0<m\le m_B$ and every $0<\mu<m$,
\begin{equation}\label{eq:strict}
  \boxed{\ I(m)<I(\mu)+I(m-\mu).\ }
\end{equation}
In particular, equality cannot occur at an interior splitting in this range. The statement does not require the existence of a minimizer for $I(s)$ at any mass $s$.
\end{theorem}

Our principal intermediate estimate, with $e(s)=I(s)/s$, is
\begin{equation}\label{eq:gap-intro}
 e(a)-e(b)\ge
 \frac{(b-a)}{a^2}
 \left(1-\frac{a+b}{2m_B}\right)[-I(a)]>0
 \quad\text{for }0<a<b,\ a+b<2m_B.
\end{equation}
This positive defect at the level of the infimum is more informative than a strict trial-function inequality, which need not remain strict after taking infima.

\section{Basic estimates and negative constrained energy}

For $u\in X$ introduce
\[
 \begin{aligned}
 M(u)&=\norm{u}_2^2,&D(u)&=\norm{\Delta u}_2^2,&
 K(u)&=\norm{\nabla u}_2^2,\\
 T(u)&=\norm{Au}_2^2,&Q(u)&=\norm{u}_4^4,&
 S(u)&=\norm{u}_6^6.
 \end{aligned}
\]
Integration by parts, or Plancherel's theorem, gives
\begin{equation}\label{eq:expand}
 T=D-2k^2K+k^4M\ge0.
\end{equation}
The embedding $H^2(\RR^2)\hookrightarrow L^p(\RR^2)$ for finite $p\ge2$ shows that $H$ is well-defined and continuous on $X$.

\begin{lemma}[Coercive lower bound]\label{lem:lower}
For each $m>0$,
\begin{equation}\label{eq:lower}
  I(m)\ge-\frac{3\beta^2}{16\gamma}\,m> -\infty.
\end{equation}
Furthermore, a fixed-mass sequence with energies bounded above is bounded in $H^2(\RR^2)$.
\end{lemma}
\begin{proof}
For $r\ge0$, square completion gives the \emph{exact} identity
\begin{equation}\label{eq:potential-square}
 -\frac{\beta}{2}r^2+\frac{\gamma}{3}r^3
 =\frac{\gamma}{3}r\left(r-\frac{3\beta}{4\gamma}\right)^2
 -\frac{3\beta^2}{16\gamma}r.
\end{equation}
Integrating with $r=\abs{u}^2$ yields
\[
 H(u)\ge T(u)-\frac{3\beta^2}{16\gamma}M(u),
\]
which proves \eqref{eq:lower}. If $M(u_n)=m$ and $H(u_n)\le C_0$, then $T(u_n)\le C_0+3\beta^2m/(16\gamma)$. Since
\[
 \norm{\Delta u_n}_2\le\norm{Au_n}_2+k^2\sqrt m,
 \qquad K(u_n)\le\sqrt{M(u_n)D(u_n)},
\]
the $H^2$ norms remain bounded.
\end{proof}

\begin{lemma}[Strict negativity at every positive mass]\label{lem:negative}
For every $m>0$, $I(m)<0$.
\end{lemma}
\begin{proof}
For $a,b>0$, set
\[
 f_{a,b}(x,y)=\frac{1}{\sqrt{\pi ab}}
 \exp\left(-\frac{x^2}{2a^2}-\frac{y^2}{2b^2}\right),
 \qquad \norm{f_{a,b}}_2=1,
\]
and, for $R>0$, consider the complex anisotropic trial
\[
 u_R(x,y)=\sqrt m\,R^{-1} f_{a,b}(x/R,y/R)e^{ikx}.
\]
It satisfies $M(u_R)=m$. Differentiating the Gaussian and integrating shows that
\begin{equation}\label{eq:gaussian-energy}
 H(u_R)=\frac{c_2}{R^2}+\frac{c_4}{R^4},
\end{equation}
where
\begin{align}
 c_2&=\frac{2mk^2}{a^2}-\frac{\beta m^2}{4\pi ab},\label{eq:c2}\\
 c_4&=m\left(\frac{3}{4a^4}+\frac{1}{2a^2b^2}
                   +\frac{3}{4b^4}\right)
       +\frac{\gamma m^3}{9\pi^2a^2b^2}>0.\label{eq:c4}
\end{align}
For any fixed $m>0$, choose $a/b>8\pi k^2/(\beta m)$, giving $c_2<0$. Then $R^2=-2c_4/c_2$ is positive and
\[
 H(u_R)=-\frac{c_2^2}{4c_4}<0.
\]
As $u_R$ has mass $m$, the constrained infimum is strictly negative.
\end{proof}

Combining the two lemmas,
\begin{equation}\label{eq:finite-negative}
 -\frac{3\beta^2}{16\gamma}m\le I(m)<0
 \qquad\text{for all }m>0.
\end{equation}

\section{Exact adjustment of minimizing sequences}

For $\lambda>0$, let
\begin{equation}\label{eq:dilation}
 u_\lambda(x)=\lambda u(\lambda x).
\end{equation}
In two dimensions, $M(u_\lambda)=M(u)$, while
\[
 D(u_\lambda)=\lambda^4 D(u),\qquad
 K(u_\lambda)=\lambda^2 K(u),\qquad
 Q(u_\lambda)=\lambda^2 Q(u),\qquad
 S(u_\lambda)=\lambda^4 S(u).
\]
Thus
\begin{equation}\label{eq:dilated-energy}
 H(u_\lambda)=\lambda^4\underbrace{\left(D+\frac\gamma3 S\right)}_{P}
 -\lambda^2\underbrace{\left(2k^2K+\frac\beta2 Q\right)}_{B}
 +k^4M.
\end{equation}
For nonzero $u\in X$, both $P>0$ and $B>0$.

\begin{lemma}[Dilation-stationary minimizing sequences]\label{lem:adjust}
For every $a>0$ there exists a minimizing sequence $(w_n)\subset X$ such that
\begin{equation}\label{eq:w-props}
 M(w_n)=a,\qquad H(w_n)\longrightarrow I(a),
\end{equation}
and each $w_n$ satisfies the exact identity
\begin{equation}\label{eq:virial}
 \boxed{\ D(w_n)-k^2K(w_n)-\frac\beta4 Q(w_n)
                   +\frac\gamma3 S(w_n)=0.\ }
\end{equation}
No attainment assumption is needed.
\end{lemma}
\begin{proof}
For a nonzero $u$, write $z=\lambda^2$. Equation \eqref{eq:dilated-energy} is the strictly convex quadratic
\[
 Pz^2-Bz+k^4M,
\]
whose unique positive minimum occurs at $z_0=B/(2P)$. Defining $w=u_{\sqrt{z_0}}$ gives
\begin{equation}\label{eq:exact-square}
 H(u_\lambda)-H(w)=P(\lambda^2-z_0)^2\ge0.
\end{equation}
Stationarity at $\lambda=\sqrt{z_0}$ reads $2P(w)=B(w)$, which is precisely \eqref{eq:virial}.

Start with any mass-$a$ minimizing sequence $(u_n)$. Apply this adjustment individually to obtain $w_n$. The dilation preserves mass and lowers or preserves energy, so
\[
 I(a)\le H(w_n)\le H(u_n)\longrightarrow I(a).
\]
This proves \eqref{eq:w-props} and \eqref{eq:virial}. Boundedness in $H^2$ also follows from Lemma~\ref{lem:lower}, although no weak or strong limit is required here.
\end{proof}

\section{Sextic-to-quartic control}

The stationarity identity produces a stronger bound than a direct application of an arbitrary-profile Sobolev estimate.

\begin{lemma}[Square completion on the stationary dilation]\label{lem:D-bound}
If $w\in X\setminus\{0\}$ satisfies \eqref{eq:virial}, then
\begin{equation}\label{eq:stationary-square}
 D(w)+T(w)+\frac{2\gamma}{3}\int_{\RR^2}
  \abs{w}^2\left(\abs{w}^2-\frac{3\beta}{8\gamma}\right)^2\dd x
 =\rho M(w),
\end{equation}
where $\rho$ is defined in \eqref{eq:mB}. In particular,
\begin{equation}\label{eq:D-rho}
 D(w)\le\rho M(w).
\end{equation}
\end{lemma}
\begin{proof}
Using \eqref{eq:expand} and \eqref{eq:virial},
\[
 \begin{aligned}
 D+T&=2D-2k^2K+k^4M\\
 &=\frac\beta2 Q-\frac{2\gamma}{3}S+k^4M.
 \end{aligned}
\]
Also, for $r=\abs w^2$,
\[
 \begin{aligned}
 \frac{2\gamma}{3}\int r\left(r-\frac{3\beta}{8\gamma}\right)^2\dd x
 &=\frac{2\gamma}{3}S-\frac\beta2 Q
   +\frac{3\beta^2}{32\gamma}M.
 \end{aligned}
\]
Adding proves \eqref{eq:stationary-square}; all three terms on the left are nonnegative, giving \eqref{eq:D-rho}.
\end{proof}

\begin{lemma}[Fourier sup-norm estimate]\label{lem:Fourier}
For every $u\in H^2(\RR^2;\CC)$ and $\eta>0$,
\begin{equation}\label{eq:Fourier-eta}
 \norm{u}_\infty^2\le\frac{D(u)+\eta M(u)}{8\sqrt\eta}.
\end{equation}
For every nonzero $u$, optimizing in $\eta$ yields
\begin{equation}\label{eq:Fourier-opt}
 \boxed{\ \norm{u}_\infty^2\le\frac14\sqrt{M(u)D(u)}.\ }
\end{equation}
\end{lemma}
\begin{proof}
Adopt the Fourier convention $\widehat u(\xi)=\int e^{-ix\cdot\xi}u(x)\dd x$. Fourier inversion, weighted Cauchy--Schwarz, and Plancherel show
\[
 \begin{aligned}
 \abs{u(x)}^2
 &\le\frac{1}{(2\pi)^4}
 \left(\int_{\RR^2}\frac{\dd\xi}{\abs\xi^4+\eta}\right)
 \left(\int_{\RR^2}(\abs\xi^4+\eta)\abs{\widehat u(\xi)}^2\dd\xi\right)\\
 &=\left[\frac{1}{(2\pi)^2}
 \int_{\RR^2}\frac{\dd\xi}{\abs\xi^4+\eta}\right]
 [D(u)+\eta M(u)].
 \end{aligned}
\]
Polar integration gives
\[
 \frac{1}{(2\pi)^2}\int_{\RR^2}
 \frac{\dd\xi}{\abs\xi^4+\eta}
 =\frac{1}{2\pi}\int_0^\infty\frac{r\dd r}{r^4+\eta}
 =\frac{1}{8\sqrt\eta}.
\]
This proves \eqref{eq:Fourier-eta}. For $u\ne0$, both $M(u)$ and $D(u)$ are positive (a square-integrable harmonic function on $\RR^2$ is zero). The choice $\eta=D(u)/M(u)$ gives \eqref{eq:Fourier-opt}. The estimate first follows for Schwartz functions and extends by density to $H^2$.
\end{proof}

\begin{proposition}[Stationary sextic/quartic ratio]\label{prop:SQ}
For every nonzero $w$ satisfying \eqref{eq:virial},
\begin{equation}\label{eq:SQ}
 \boxed{\ \frac{S(w)}{Q(w)}\le\norm{w}_\infty^2
       \le\frac\sigma4 M(w).\ }
\end{equation}
\end{proposition}
\begin{proof}
By \eqref{eq:Fourier-opt} and \eqref{eq:D-rho},
\[
 \norm{w}_\infty^2\le\frac14\sqrt{M(w)D(w)}
 \le\frac{\sqrt\rho}{4}M(w).
\]
Also $S(w)=\int\abs{w}^2\abs{w}^4\dd x\le\norm w_\infty^2Q(w)$; since $w\ne0$, $Q(w)>0$.
\end{proof}

\section{A quantitative energy-per-mass gap}

\begin{proposition}[Strict decrease at the infimum]\label{prop:gap}
Let $0<a<b$ and assume $a+b<2m_B$. Set
\[
 e(s)=\frac{I(s)}s,\qquad
 h=1-\frac{a+b}{2m_B}>0.
\]
Then
\begin{equation}\label{eq:gap}
 \boxed{\ e(a)-e(b)\ge\frac{(b-a)h}{a^2}[-I(a)]>0.\ }
\end{equation}
\end{proposition}
\begin{proof}
Choose any minimizing sequence $(w_n)$ at mass $a$, adjusted as in Lemma~\ref{lem:adjust}. Set $t=b/a>1$. Amplitude rescaling has mass $M(\sqrt t\,w_n)=ta=b$ and transforms the energy according to
\begin{equation}\label{eq:amp}
 \frac{H(\sqrt t\,w_n)}t-H(w_n)
 =(t-1)\left[-\frac\beta2 Q(w_n)
       +\frac{\gamma(t+1)}3S(w_n)\right].
\end{equation}
By Proposition~\ref{prop:SQ}, using $M(w_n)=a$,
\[
 \begin{aligned}
 -\frac\beta2 Q(w_n)+\frac{\gamma(t+1)}3S(w_n)
 &\le\left[-\frac\beta2+\frac{\gamma\sigma(t+1)a}{12}\right]Q(w_n)\\
 &=-\frac\beta2\left(1-\frac{a+b}{2m_B}\right)Q(w_n)\\
 &=-\frac{\beta h}{2}Q(w_n).
 \end{aligned}
\]
Since $T\ge0$ and $S\ge0$, we have
\[
 H(w_n)\ge-\frac\beta2Q(w_n).
\]
Consequently,
\begin{equation}\label{eq:uniform-inf}
 \begin{aligned}
 \frac{I(b)}t
 &\le\frac{H(\sqrt t\,w_n)}t\\
 &\le H(w_n)-(t-1)\frac{\beta h}{2}Q(w_n)\\
 &\le[1+(t-1)h]H(w_n).
 \end{aligned}
\end{equation}
Pass to the limit $n\to\infty$ and divide by $a$:
\[
 e(b)\le[1+(t-1)h]e(a).
\]
It follows that
\[
 e(a)-e(b)\ge-(t-1)he(a)
 =\frac{(b-a)h}{a^2}[-I(a)].
\]
The right side is strictly positive by Lemma~\ref{lem:negative}. In particular the positive defect is independent of the minimizing-sequence index and cannot vanish upon taking the infimum.
\end{proof}

\begin{remark}[The endpoint and larger-mass pairs]
If $0<a<b\le m_B$, then $a+b<2m_B$, including when $b=m_B$. The argument also applies to some comparisons with $b>m_B$, provided $a+b<2m_B$; it does not assert monotonicity for arbitrary masses.
\end{remark}

\section{Strict binding and the frozen-parameter corollary}

\begin{proof}[Proof of Theorem~\ref{thm:main}]
Fix $0<m\le m_B$ and $0<\mu<m$. Since
\[
 \mu+m<2m\le2m_B,
 \qquad (m-\mu)+m<2m\le2m_B,
\]
Proposition~\ref{prop:gap}, applied to $(\mu,m)$ and $(m-\mu,m)$, yields
\[
 e(\mu)>e(m),\qquad e(m-\mu)>e(m).
\]
Both $\mu$ and $m-\mu$ are strictly positive; multiplying the two inequalities by those masses and adding gives
\[
 \begin{aligned}
 I(\mu)+I(m-\mu)
 &=\mu e(\mu)+(m-\mu)e(m-\mu)\\
 &>[\mu+(m-\mu)]e(m)=I(m).
 \end{aligned}
\]
This is \eqref{eq:strict} for every interior split.
\end{proof}

\begin{corollary}[Exact frozen-parameter check]\label{cor:frozen}
Consider $k=1$, $\beta=\gamma=b_0$, and $m=m_0$, where the parameters are the exactly represented binary64 values
\[
 b_0=\texttt{0x1.05f1d178c75ebp+2},\qquad
 m_0=\texttt{0x1.45b8a0bad2714p+1}.
\]
Viewed as rational numbers, they equal
\[
 b_0=\frac{4608178181141995}{1125899906842624},\qquad
 m_0=\frac{1432536998431173}{562949953421312}.
\]
For these coefficients, $m_0<m_B$. Therefore the strict-binding inequality holds at mass $m_0$ for all $0<\mu<m_0$.
\end{corollary}
\begin{proof}
Because $k=1$ and $\beta=\gamma=b_0$,
\[
 \rho=1+\frac{3b_0}{32}
 =\frac{49853331562389953}{36028797018963968},
 \qquad m_B=\frac{3}{\sqrt\rho}.
\]
Exact rational arithmetic (before rounding any square roots) gives
\begin{equation}\label{eq:frozen-rational}
 9-\rho m_0^2=
 \frac{454708712324087584212476662678391978900775591}
 {11417981541647679048466287755595961091061972992}>0.
\end{equation}
As all quantities are positive, this proves $m_0<3/\sqrt\rho=m_B$. For orientation only, the decimal values are
\[
 m_0=2.54469689485712\ldots,\qquad
 m_B=2.550345636351549\ldots.
\]
Apply Theorem~\ref{thm:main}.
\end{proof}

\begin{remark}[Scope of the frozen corollary]
The notation ``D4'' used in the associated computational project refers there to a particular square-symmetric configuration in two spatial dimensions. Corollary~\ref{cor:frozen} is a statement about an unrestricted whole-space functional on $\RR^2$; it does not establish identification with a computed periodic branch, uniqueness of minimizers, spectral stability of a branch, or a lifetime classification for four spatial dimensions.
\end{remark}

\section{Discussion and limitations}

The proof uses two complementary scalings. The mass-preserving spatial dilation $u\mapsto\lambda u(\lambda\cdot)$ supplies an exact identity on adjusted minimizing sequences. The amplitude scaling $u\mapsto\sqrt t\,u$ then compares two masses. The positive sextic amplitude contribution prevents a universal comparison for \emph{all} profiles; the stationary ratio estimate \eqref{eq:SQ} controls that contribution on the adjusted sequences. In particular, negativity of $I(m)$ alone would not suffice to justify the comparison.

The threshold $m_B$ is a \emph{sufficient} mass bound, not a proved sharp boundary of strict binding. This paper does not claim equality, failure of binding, or optimality of its constants for masses above $m_B$. Likewise, strict binding is distinct from attainment and compactness: concentration--compactness can use strict subadditivity to exclude certain splitting scenarios, but any existence or dynamical stability theorem requires additional hypotheses and arguments.

For completeness, ordinary subadditivity $I(a+b)\le I(a)+I(b)$ follows by approximating mass-$a$ and mass-$b$ near-minimizers by compactly supported smooth profiles, restoring their exact masses by a factor tending to one, and translating one profile until the supports are separated. The mass, quadratic derivative form, and local nonlinear terms then add exactly; approximation and continuity of $H$ allow passage to the infima. The strict conclusion of Theorem~\ref{thm:main} is not a consequence of that elementary construction, but of Proposition~\ref{prop:gap}.

\paragraph{Reproducibility.}
The associated development notes, including independent dilation and arithmetic audits, are available in the public repository at
\begin{center}
\small\url{https://github.com/rickyjreyes/geometry_of_resonance/tree/main/docs/d4_binding}.
\end{center}
The theorem and its proof above are self-contained and do not depend on numerical evidence or external computational artifacts.

\begin{thebibliography}{9}
\bibitem{FJMM2022}
A.~J. Fern\'andez, L. Jeanjean, R. Mandel, and M. Mari\c{s},
\emph{Non-homogeneous Gagliardo--Nirenberg inequalities in $\RR^N$ and application to a biharmonic non-linear Schr\"odinger equation},
Journal of Differential Equations \textbf{330} (2022), 1--65.
\href{https://doi.org/10.1016/j.jde.2022.04.037}{doi:10.1016/j.jde.2022.04.037}.

\bibitem{MCF2024}
Z. Ma, X. Chang, and Z. Feng,
\emph{Normalized ground state of a mixed dispersion nonlinear Schr\"odinger equation with combined power-type nonlinearities},
Electronic Journal of Differential Equations \textbf{2024} (2024), no.~29, 1--20.
\href{https://doi.org/10.58997/ejde.2024.29}{doi:10.58997/ejde.2024.29}.
\end{thebibliography}

\end{document}